\chapter{Inference under the PCH Model}

%%%%%%%%%%%%%%%%%%%%%%%%%%%%%%%%%%%%%%%%%%%%%%%%%
% Introduction
%%%%%%%%%%%%%%%%%%%%%%%%%%%%%%%%%%%%%%%%%%%%%%%%%
\section{Introduction}

In phylogenetics, we have a collection of objects whose ancestral history we hope to infer
using only modern-day information. This shared history is summarized by a \textit{phylogenetic tree},
a graphical tree whose leaves are labeled by these objects. Classically, the objects were biological species,
the data were DNA sequences, and the phylogenetic tree encoded the sequence of
speciation events that occurred over time to produce the species we see today. Nonetheless, the tools of phylogenetics are far
more broadly applicable: they have been used in medicine to better understand the
progression of cancer~\cite{somarelli2017_phylooncology_cancer_phylogenetic_analysis}, in astrophysics as a hierarchical clustering procedure for astral
 bodies~\cite{fraixburnet2015_multivariate_classification_astronomy}, in terrane analysis to study the geological history of specific
 regions~\cite{michaux2018_cladistic_methods_terrane_analysis}, and in many other academic disciplines.

In this chapter, we study a specific application to linguistics, where the objects of interest are different
languages and the underlying data are \textit{linguistic characters}. These characters encode features
that vary across languages and are commonly divided into three broad categories: \textit{phonological},
 \textit{morphological}, and \textit{lexical}.

 Phonological characters describe sound patterns,
 such as whether a language distinguishes tones or permits consonant clusters. Morphological characters
 describe grammatical structure, such as how past tense or plurality is expressed. Lexical characters describe the word
 forms used to express specific ideas, often by grouping languages according to whether their words for
  common concepts such as ``fire,'' or ``mother'' descend from a common historical source. 
  
  Lexical characters are
  particularly interesting as they can be multi-state---several different words can simultaneously share
  the same meaning---a property known as \textit{polymorphism}. For example, in English the words ``stone'' and ``rock'' fill
  the same semantic slot. We are specifically interested in the model introduced in
\cite{canby2024_polymorphism_linguistic_phylogenetics}, which explicitly models this polymorphism. We refer to this model as the
\textit{Polymorphic CHaracter (PCH) model}, in a similar vein to the follow-up paper \cite{evans2026_canby_identifiability}.

%%%%%%%%%%%%%%%%%%%%
% PCH Model
%%%%%%%%%%%%%%%%%%%%
\subsection{The PCH Model}

The key feature of the PCH model is that it allows for polymorphism: a given character $c$ may take on multiple different
states simultaneously. To differentiate between the set-valued collection of states a character currently takes on and the individual
states themselves, we call the former the \textit{macro-state} of the character and the latter its individual \textit{micro-states}. For
example, if the macro-state of a character $c$ at a given point of time is $\{\alpha, \beta\}$ then it currently exhibits two micro-states: $\alpha$ and  $\beta$.
Under the PCH model, it is assumed that micro-states may be born or die as the character evolves along the tree, so that the current
\textit{number} of micro-states present in the current macro-state evolves according to a standard birth-death process, which we refer to as the \textit{cardinality process}.

More formally, suppose that $\cT^*$ is the underlying rooted phylogenetic tree of languages. Given some collection of characters $\cC$,
associated to each character $c$ is a set $\Sigma^{(c)}$
of possible states, birth/death rates $\{\lambda_i^{(c)}\}, \{\mu_i^{(c)}\}$ controlling the cardinality process, a probability measure $\nu^{(c)}$
on $\Sigma^{(c)}$ for selecting new states, and a probability measure $\rho^{(c)}$ on the set ${\Sigma^{(c)} \choose *}$ of nonempty finite subsets of $\Sigma^{(c)}$ for
determining the macro-state of the root.

More specifically, let $[r, \ell]$ denote the unique path from the root $r$ of $\cT^*$ to some leaf $\ell \in \cT^*$. For $x \in [r, \ell]$,
macro-state $c(x)$ evolves down this path as follows:
%
\begin{enumerate}[label=(\roman*)]
    \item The macro-state $c(r)$ of the root is sampled from $\rho^{(c)}$.
    \item If $|c(x)| = i$, a new micro-state is born at rate $\lambda_i^{(c)}$ and an existing micro-state dies at rate $\mu_i^{(c)}$ respectively.
    \item If a death occurs, a micro-state $\alpha \in c(x)$ is chosen uniformly at random to delete. If a birth occurs and $\nu^{(c)}(\Sigma^{(c)} \setminus c(x)) > 0$, a new micro-state $\alpha \notin c(x)$ is chosen proportional to $\nu^{(c)}$, namely via
                \[\nu_x^{(c)}(A) := \frac{\nu^{(c)}(A \setminus c(x))}{\nu^{(c)}(\Sigma^{(c)} \setminus c(x))}.\]
\end{enumerate}
%
When a branch-point $v$ of $\cT^*$ is reached, conditional on $c(v)$ the macro-state evolves independently down each child branch starting from
state $c(v)$.

A micro-state $\alpha \in \Sigma^{(c)}$ with $\nu^{(c)}(\{\alpha\}) > 0$ is called \textit{homoplastic}. These states
can appear independently in different lineages, rather than from shared ancestry. We denote the set of homoplastic states
and non-homoplastic states by $\Sigma_H^{(c)}, \Sigma_N^{(c)}$ respectively. Homoplastic states in a sense complicate the
inference process, as two languages can share the homoplastic trait without necessarily having any recent shared ancestry.

Throughout this chapter, we assume the cardinality process has \textit{immigration-death rates}
%
\begin{equation}
    \lambda_i^{(c)} = \lambda^{(c)}, \quad \mu_i^{(c)} = i \mu^{(c)} 1\{i > 1\},
\end{equation}
%
\noindent for $i \geq 1$. Namely, the birth rate is constant, and the death rate is proportional to the number of states,
except when only one state is present. These assumptions ensure each character
has at least one micro-state at any point in $\cT^*$ and that the number of micro-states
cannot blow-up in finite time.

Since the upcoming inference procedure is based solely on non-homoplastic states, we also make the
following assumptions about $\nu^{(c)}$ and $\rho^{(c)}$.

\begin{assumption}[1]
    Non-homoplastic states occur sufficiently often, specifically:
    %
    \begin{enumerate}[label=(\roman*)]
        \item The root macro-state contains a non-homoplastic micro-state with positive probability
              %
                \[\rho^{(c)} \left(\left\{A \in {\Sigma^{(c)} \choose *} \; : \; A \cap \Sigma^{(c)}_N \neq \emptyset \right\}\right) > 0.\]
              %
        \item There is positive probability of adding a non-homoplastic micro-state to the current macro-state when a birth occurs
                %
                \[\nu^{(c)}(\Sigma_N^{(c)}) > 0.\]
                %
    \end{enumerate}
\end{assumption}

The full PCH model additionally incorporates another mechanism of character transmission known
as \textit{borrowing} which allows characters to transfer laterally throughout the tree, analogous to
horizontal gene transfer in biological contexts. In this chapter, we only consider the case where no borrowing
occurs, but interested readers can find further details in \cite{canby2024_polymorphism_linguistic_phylogenetics}.


%%%%%%%%%%%%%%%%%%%%%%%%%
% ASTRAL
%%%%%%%%%%%%%%%%%%%%%%%%%
\subsection{ASTRAL}

To infer the underlying tree from many independently evolving characters, we adapt ideas from
quartet-based species-tree estimation. In phylogenetic inference, one often wants to combine information from many sources, such as genetic information across
a collection of genes or, in our case, linguistic characters. In \textit{consensus methods}, each source is
used to generate an estimate $g_i$ of the true phylogenetic tree $\cT^*$, and then these estimates $\{g_i\}_{i=1}^n$
are aggregated. In the bioinformatics literature, the trees $g_i$ are called the \textit{gene trees}, while $\cT^*$ is
referred to as the \textit{species tree}. We will occasionally utilize this vocabulary.

ASTRAL~\cite{mirarab2014_astral_species_tree_estimation} is one commonly used aggregation procedure. It scores candidate
phylogenetic trees $\cT$ based on the number of \textit{quartets}, four-leaf induced subtrees, they share with the estimated trees $\{g_i\}$. Specifically,
it solves the optimization problem
%
\begin{equation}
    \hat{\cT} = \argmax_\cT  \sum_{q \in \cQ(\cT)} W(q),
    \label{eqn:astral_quartet_support_optimization_problem}
\end{equation}
%
\noindent where $W(q)$ is some measure of the support for the quartet $q$ and 
$\cQ(\cT)$ is the set of quartets present in the species-tree $\cT$. Commonly,
%
    \[W(q) = \sum_{i=1}^n 1 \{q \in g_i\},\]
%
\noindent the number of gene trees in which the quartet appears. Later on, we will introduce
other measures of the quartet support specific to the PCH model.

Since the set of phylogenetic trees is super-exponentially large, it is intractable to solve \eqref{eqn:astral_quartet_support_optimization_problem}
exactly for trees with more than 20 or so leaves. For larger trees, a
restricted search space is required, and one possible strategy is to limit the bipartitions of the leaves the candidate trees $\cT$ can induce.

A \textit{bipartition} of a set $S$ is a partition of $S = A \cup B$ into two nonempty, disjoint subsets.
Every edge $e \in E(\cT)$ of a phylogenetic tree $\cT$ induces a bipartition of the leaves by considering the two connected components
of the graph $(V(\cT), E(\cT) \setminus \{e\})$. Moreover, Buneman's splits-equivalence theorem (e.g. \cite[Theorem~3.1.4]{semple2003_phylogenetics}) shows $\cT$
is completely characterized by its set $\cX(\cT)$ of bipartitions. Since each $g_i$ is a noisy estimate of $\cT^*$, we
might expect that $\cX(\cT^*) \subseteq \cX_G := \bigcup_i \cX(g_i)$, suggesting
the restricted search space
%
\[\Big\{\cT \; : \; \cX(\cT) \subseteq \cX_G\Big\}.\]
%
\noindent More generally, given any set $\cX$ of bipartitions, we can solve
%
%
\begin{equation}
    \hat{\cT} = \argmax_{\cX(\cT) \subseteq \cX}  \sum_{q \in \cQ(\cT)} W(q).
    \label{eqn:astral_heuristic_optimization_problem}
\end{equation}
%
%
Using a clever dynamic-programming principle, ASTRAL efficiently solves this bipartition-constrained
optimization problem. For polynomially-sized $\cX$, the optimization problem~\eqref{eqn:astral_heuristic_optimization_problem}
can be feasibly solved for trees with thousands of leaves. Moreover, when the trees $g_i$ are generated under the
multispecies coalescent (MSC) model~\cite{rannala2020_multispecies_coalescent_species_tree_inference}, \cite{hill2020_species_tree_joint_coalescence_duplication}, its statistical properties
are well-understood. In particular, choosing $\cX$ to contain $\cX_G$ is asymptotically sufficient because,
as the number of genes grows, $\cX_G$ contains every bipartition of $\cT^*$ with probability tending to one.

The paper~\cite{shekhar2018_how_many_genes_are_enough} showed for the exact version~\eqref{eqn:astral_quartet_support_optimization_problem}
that $n = \Theta(f^{-2} \log(k))$ independent gene trees are required to recover $\cT^*$ with high probability under the MSC model, where $f$ is the minimum
branch length in the tree. It has been empirically observed that the relaxed version~\eqref{eqn:astral_heuristic_optimization_problem} requires
significantly more data to obtain similar recovery rates~\cite{shekhar2018_how_many_genes_are_enough}. Thus there is a trade-off
between computational and sample efficiency.

As discussed in \cite{uricchio2016_bipartition_cover}, \cite{mcnulty2026_improved_bipartition_cover_bound}, the main shortcoming of the relaxed problem
is that it is quite unlikely that $\cX(\cT^*) \subseteq \cX_G$, since the bipartitions induced by edges near the root are very
difficult to recover. Several heuristics have been proposed to extend the set $\cX_G$ in an effort to improve its coverage
properties and the empirical performance of the relaxed problem~\cite{mirarab2015_astral_ii_species_tree_estimation}, \cite{zhang2018_astral_iii_species_tree_reconstruction}, \cite{dibaeinia2021_fastral_scalable_phylogenomics}.
Alternatively \cite{zhang2022_weighting_gene_tree_uncertainty} proposes a weighted version of~\eqref{eqn:astral_quartet_support_optimization_problem}
which incorporates quartet-support information from the gene-tree estimation stage, which down-weights possibly misleading quartets.

ASTRAL recovers only the underlying unrooted tree. When an appropriate outgroup is available, it can be used to root the inferred tree.

%%%%%%%%%%%%%%%%%%%%%%%%%%%%%%%%%%%%
% ASTRAL-Based Inference for PCH Model
%%%%%%%%%%%%%%%%%%%%%%%%%%%%%%%%%%%%
\subsection{PCH-ASTRAL: Phylogenetic Inference in the PCH Model}

For data generated under the PCH model, quartet support can be defined directly from the
shared non-homoplastic states of the observed languages. Let $q = uv||xy$ denote the unrooted
quartet topology that groups $u$ with $v$ and $x$ with $y$.

\begin{center}
\begin{tikzpicture}[scale=0.9, every node/.style={font=\small}]
    \coordinate (a) at (-1, 0);
    \coordinate (b) at (1, 0);
    \coordinate (u) at (-2.5, 0.9);
    \coordinate (v) at (-2.5, -0.9);
    \coordinate (x) at (2.5, 0.9);
    \coordinate (y) at (2.5, -0.9);

    \draw (u) -- (a) -- (b) -- (x);
    \draw (v) -- (a);
    \draw (y) -- (b);

    \node[left] at (u) {$u$};
    \node[left] at (v) {$v$};
    \node[right] at (x) {$x$};
    \node[right] at (y) {$y$};
\end{tikzpicture}
\end{center}

\noindent Given a character $c$, define the support $M^{(c)}_q$ of the quartet $q = uv||xy$ via
%
\begin{equation*}
    M^{(c)}_{q}:= \Big\{(\alpha, \beta) \in \Sigma_N \times \Sigma_N \; : \; \alpha \neq \beta, \quad \alpha \in U \cap V \cap \bar{X} \cap \bar{Y}, \quad \beta \in \bar{U} \cap \bar{V} \cap X \cap Y\Big\},
\end{equation*}
%
where $U := c(u)$, and so on, with $\bar U:=\Sigma^{(c)}\setminus U$, and similarly for the other leaves. Namely, $M_q$ consists of the pairs of non-homoplastic states that are seemingly ``split'' across the quartet $q$.
For the true quartet topology, \cite[Theorem 1]{evans2026_canby_identifiability} shows this support is larger in expectation than the support assigned to
the two alternatives:
%
\begin{theorem}
\label{thm:correct_topo_max_expected_support}
    Let $s=\{u,v,x,y\}$. If $\cT_s^*=uv||xy$, then
    \[\E |M_{uv || xy}| > \E |M_{ux || vy}| = \E |M_{uy || vx}|,\]
    %
    \noindent so the correct topology has the highest expected support.
\end{theorem}
%
\noindent For each quartet $q$, they define its (empirical) support via
%
\begin{equation}
    S_q := \frac{1}{n} \sum_{i=1}^n |M_q^{(c_i)}|.
    \label{eqn:quartet_support_averaged}
\end{equation}
%
\noindent Moreover, for each $s = \{u,v,x,y\} \in {[k] \choose 4}$ they then let
%
\begin{equation}
    \hat{q}_s := \argmax_{q \in A_s} S_q,
    \label{eqn:astral_based_canby_inference_original}
\end{equation}
%
\noindent with ties broken uniformly, where $A_s$ is the set of three quartets with leaf-set $s$. Namely, $\hat{q}_s$ is the
quartet topology with maximal empirical support.
By the law of large numbers and Theorem~\ref{thm:correct_topo_max_expected_support}, the correct
quartet topology becomes overwhelmingly likely for large $n$. They then run ASTRAL on the quartet set $\{\hat{q}_s\}$
to generate an estimate $\hat{\cT}$ of $\cT^*$. The authors of \cite{evans2026_canby_identifiability} call this algorithm \textit{PCH-ASTRAL}, and
additionally propose a weighted version, although we will
only focus on the unweighted version described above.


%%%%%%%%%%%%
% Summary of contributions
%%%%%%%%%%%%
\subsection{Summary of Contributions}

In Section~\ref{sec:basic_results}, we construct a coupling of the PCH model to a simple immigration-death
process which allows for many direct computations. This will be the basis of most of our rigorous analysis.
Using this, we develop some basic properties of both the original PCH model and of this coupled process.


Afterward, in Section~\ref{sec:bipartition_covers_canby}, we show how the search space for PCH-ASTRAL can be
significantly restricted while maintaining statistical consistency. This produces a much smaller candidate set
and is intended to make the underlying ASTRAL procedure more scalable. While the
restriction increases the algorithm's \textit{sample complexity}, the number of data samples required to recover
$\cT^*$ with a prescribed probability of success, we further introduce two heuristics intended to mitigate this
loss in statistical efficiency.

From there, in Section~\ref{sec:pch_astral_sample_complexity}, we show under both the PCH model and its immigration-death approximation that
PCH-ASTRAL requires roughly
%
    \[O\left(\kappa^2L^{-2}\log(k)+\kappa L^{-1}\log^2(k)\right)\]
%
\noindent characters to recover $\cT^*$ with high probability, where
$\kappa$ is the worst-case $\Psi_{1/2}$ pseudo-norm of the quartet supports and
%
    \[L := \min_{s \in {[k] \choose 4}} \min_{q \in A_s \setminus \{\cT_s^*\}} \E \left[|M_{\cT_s^*}| - |M_q|\right]\]
%
is the minimal margin the correct quartet topology has over incorrect ones. In Lemma~\ref{lem:canby_expected_margin}, we
give a lower bound for $L$ that scales at least quadratically in the minimum branch length $f$
when the other geometric parameters are fixed.

Then, Section~\ref{sec:canby_alternative_astral_procedure} studies an alternative character-voting
procedure. This procedure, which we call \textit{PCH-ASTRAL(sc)}, keeps character-level quartet information
separate by converting each informative character-quartet pair into a single topology vote, rather than averaging
 support magnitudes across characters. When a certain \emph{anomaly-free} property holds, we show PCH-ASTRAL(sc) requires
 %
    \[O(p_{\min}^{-1}\Delta^{-2}\log(k))\]
 %
 characters, where
%
    \[\Delta := \min_{s \in {[k] \choose 4}} \min_{q \in A_s \setminus \{\cT_s^*\}} \left[\P_s(\cT_s^*) - \P_s(q)\right], \quad \P_s(q) := \P\left(q = \argmax_{q' \in A_s} |M_{q'}| \; \Big\vert \; \max_{q' \in A_s} |M_{q'}| > 0\right).\]
%
again gives some notion of the margin and
%
    \[p_{\min} := \min_{s \in {[k] \choose 4}} \P\left(\max_{q \in A_s} |M_q^{(c)}| > 0\right),\]
%
\noindent is the worst-case probability a specific leaf-set has \emph{some}, possibly incorrect, topology
with non-zero support. In Lemma~\ref{lem:low_polymorphism_anomaly_free}, we show this anomaly-free property holds
under the ID approximation in the low-polymorphism regime. Because the supports $|M_q|$
are moderately heavy-tailed, PCH-ASTRAL(sc) may be more robust to influential characters.




%%%%%%%%%%%%
% Notation
%%%%%%%%%%%%
\subsection{Notation}

Let $\cT^*$ be the true species tree, with root $r$ and leaf set $\cL(\cT^*)$,
which we identify with $[k]$. For $i,j\in\cT^*$, let $[i,j]$ be their unique
connecting path and $\ell_{ij}$ its length. For fixed $\mu, \lambda$, which
will be clear from context, we write
%
\[\eta_T:=e^{-\mu T},\qquad \eta_{ij}:=e^{-\mu\ell_{ij}},\qquad
m_\theta(T):=\theta e^{-\mu T}+\frac{\lambda}{\mu}(1-e^{-\mu T}).\]
%
For a four-leaf set $s=\{u,v,x,y\}$, write $q=uv||xy$ for a particular
unrooted topology and
%
\[A_s:=\{uv||xy,ux||vy,uy||vx\}\]
%
for all three possibilities. The topology induced by $\cT^*$ is $\cT_s^*$.
For any tree $\cT$, let $\cQ(\cT)$ be its set of induced quartets; thus
$|\cQ(\cT)|={k\choose4}$ when $\cT$ has $k$ leaves.

Unless stated otherwise, characters are assumed iid and any referenced micro-states are assumed
non-homoplastic. We occasionally abbreviate macro-states by upper-case letters,
for example $I=c(i)$, $J=c(j)$, and $K=c(k)$. Lastly, we say $U$ \emph{first-order stochastically
dominates} $V$ if $\E h(U)\geq\E h(V)$ for every bounded increasing function $h$, and
denote this by $U\succeq_{st}V$.

%%%%%%%%%%%%%%%%%%%%%%%%%%%%%%%%%%%%%%%%%%%%%%%%%%%%%%%%%%%%
% Basic Results
%%%%%%%%%%%%%%%%%%%%%%%%%%%%%%%%%%%%%%%%%%%%%%%%%%%%%%%%%%%%%
\section{Basic Results}
\label{sec:basic_results}

%%%%%%%%%%%%%%%%%%%%%
% Coupling to ID Process
%%%%%%%%%%%%%%%%%%%%%
\subsection{Coupling to Immigration-Death Processes}

The constraint $\mu_1 = 0$ and the renormalization by $\nu^{(c)} (\Sigma^{(c)} \setminus c(x))$ make 
the PCH model hard to analyze as they add a certain
dependence structure between the different micro-states of a specific character, and they make
the birth-rate of non-homoplastic states time inhomogeneous. If we drop these conditions, we
get a classical immigration-death (ID) process, which is much more tractable. This approximation
is reasonable in the parameter regime where the number of micro-states in a given branch is rarely one, which occurs
when $\lambda$ is sufficiently large relative to $\mu$, and when homoplasy is fairly rare, so renormalization has little
influence.

Our key tool will be a coupling of the PCH model to a standard ID process
 run along the branches of our species tree $\cT^*$, and whose initial population 
 size is Poissonian. It relies on two simple observations
 about how the set of nonhomoplastic states evolves over time:

\begin{itemize}
    \item Births always occur at rate at least $\lambda_N := \lambda \nu(\Sigma_N)$ and at most $\lambda$.

    \item Deaths always occur at rate at most $\mu$ per individual.
\end{itemize}
%
\noindent In particular, we can couple an ID process $(Y_t)_{t \in \cT^*}$ to our PCH model $(X_t)_{t \in \cT^*}$ so
%
\[ Y_t \subseteq X_t \cap \Sigma_N, \quad \forall t \in \cT^*.\]
%
\noindent In this way, we only reduce the number of nonhomoplastic states we observe in the leaves, which heuristically makes inference more challenging.

To produce this coupling, we start by coupling the root states. Suppose
$|X_0 \cap \Sigma_N| \succeq_{st} Z_\theta$ for some
$Z_\theta \sim \mathrm{Poisson}(\theta)$ and $\theta \geq 0$. Since we can
always choose $\theta = 0$, such a $\theta$ always exists. Then there exists a coupling so $|X_0 \cap \Sigma_N| \geq Z_\theta$ 
almost surely. Let $Y_0$ be a uniformly at random chosen subset of size $Z_\theta$ from $X_0 \cap \Sigma_N$.

From here, we use rate-$\lambda$ exponential clocks along the branches of
$\cT^*$ to determine possible births. If a clock rings at location $t \in \cT^*$,
sample a new state $\alpha \notin X_{t^-}$ for the PCH model as usual, namely
according to the conditional distribution
%
\[\frac{\nu(\cdot \cap (\Sigma \setminus X_{t^-}))}{\nu( \Sigma \setminus X_{t^-})}.\]
%
\noindent Let $X_t = X_{t^-} \cup \{\alpha\}$, and update our ID process $Y_t$ as
 follows: if $\alpha \notin \Sigma_N$, then let $Y_t = Y_{t^-}$. Namely, ignore
 the homoplastic states. Otherwise, let
%
\[Y_t = \begin{cases}
    Y_{t^-} \cup \{\alpha\} & \mathrm{w.p.}  \quad \nu(\Sigma \setminus X_{t^-}),\\
    Y_{t^-} & \mathrm{w.p.} \quad 1-\nu(\Sigma \setminus X_{t^-}). \\
\end{cases}\]
%
\noindent This ensures that the birth rate in $(Y_t)$ along a given branch
is always $\lambda_N$. Similarly, assign each state $\alpha$ independent
rate-$\mu$ exponential clocks along the branches of $\cT^*$. If a clock for
$\alpha$ rings at location $t \in \cT^*$, let
%
\[X_t = \begin{cases}
    X_{t^-} \setminus \{\alpha\} & X_{t^-} \neq \{\alpha\},\\
    \{\alpha\} & \mathrm{otherwise}.
\end{cases}\]
%
\noindent Then let $Y_t = Y_{t^-} \setminus \{\alpha\}$. Namely, $(X_t)$ simply ignores the deaths causing the state set to become empty, while $(Y_t)$ does not.

Heuristically, this coupling well-approximates the non-homoplastic states of the PCH model in the
high-polymorphism regime where it is rare for there to be just a
single state, and when there is little homoplasy so the renormalization of
$\nu$ does not dramatically affect the birth rate of non-homoplastic states
over time.

To this end, let $ID(\theta, \lambda, \mu; \cT^*)$ denote an immigration-death process on the species tree $\cT^*$ with birth rate 
$\lambda$, death rate $\mu$, and with an initial population of size $\mathrm{Poisson}(\theta)$. The final assumption 
ensures that the number of states present at any point of $\cT^*$ is Poisson distributed, making the process particularly simple to analyze.

To start, we provide an estimate of how close this coupled process
is to the original PCH model, and quantify the main factors in any discrepancies.
For $t \in [r, \ell]$, let
%
\[
    D_t:=\left|(X_t\cap\Sigma_N)\setminus Y_t\right|,
    \qquad \lambda_H:=\lambda - \lambda_N.
\]
%
\noindent count the number of non-homoplastic states on which the coupled process and original PCH
model disagree. Then

\begin{proposition}[Coupling discrepancy]
\label{prop:long_term_id_coupling_discrepancy} If $\E D_0 < \infty$, then 
%
\[
    \E D_t
    \leq e^{-\mu t}\E D_0
    +\frac{\lambda_H}{\mu}(1-e^{-\mu t})
    +\mu\int_0^t e^{-\mu(t-u)}\P(|X_u|=1)\,du.
\]
%
Moreover,
%
\[
    \limsup_{t\to\infty}\E D_t
    \leq
    \frac{\lambda_H}{\mu}
    +\frac{\lambda/\mu}{e^{\lambda/\mu}-1}.
\]
%
\end{proposition}

\begin{proof}

Note $D_t \mapsto D_t + 1$ when either a non-homoplastic birth occurs in $X_t$ and is ignored in $Y_t$ 
or when $X_t = Y_t = \{\alpha\}$ for $\alpha \in \Sigma_N$, and a death clock rings. This occurs at rate
%
    \[\frac{\lambda_N \nu(X_t)}{1-\nu(X_t)} + \mu 1\{|X_t| = |Y_t| = 1\}.\]
%
\noindent Similarly, we see $D_t \mapsto D_t - 1$ when one of the states in $(X_t \cap \Sigma_N) \setminus Y_t$ dies,
at rate
%
    \[\mu D_t 1\{|X_t| > 1\}.\]
%
\noindent Thus we see
%
    \[\frac{d}{dt} \E D_t = \lambda_N \E \left(\frac{\nu(X_t)}{1-\nu(X_t)}\right) - \mu \E D_t + \mu \P(|X_t| = 1).\]
%
Since $\nu$ is diffuse on $\Sigma_N$, $\nu(X_t)\leq1-\nu(\Sigma_N)$. Thus
%
\[
    \frac{d}{dt}\E D_t
    \leq \lambda_H+\mu\P(|X_t|=1)-\mu\E D_t.
\]
%
Gr\"onwall's inequality gives the finite-time bound. By
Lemma~\ref{lem:stationary_dist_state_cardinality_process},
%
\[
    \P(|X_t|=1)\longrightarrow
    \frac{\lambda/\mu}{e^{\lambda/\mu}-1},
\]
%
so the bounded convergence theorem gives the asymptotic bound.
\end{proof}

Thus the coupled process is close in expected population discrepancy when
homoplasy is rare and $\lambda/\mu$ is large. In the low-polymorphism regime,
the singleton term is instead close to one, reflecting the fact that the ID
process may lose its final state while the PCH process cannot.

%%%%%%%%%%%%%%%%%%%%%
% properties of PCH Model
%%%%%%%%%%%%%%%%%%%%%
\subsection{Properties of PCH Model}

In this section, we derive some basic properties of the PCH model. The derivations often rely on the
coupling we introduced in the previous section.

Given a path $[r, \ell]$ from the root $r$ of $\cT^*$ to some leaf $\ell$, let $S_t = |c(t)|$ for $t \in \cT^*$ be the cardinality
process of the PCH model. This is a simple birth-death process with birth rates
$\lambda_i \equiv \lambda$ and death rates $\mu_i = i\mu 1\{i > 1\}$. Hence

% stationary distribution
\begin{lemma}\label{lem:stationary_dist_state_cardinality_process}
    The stationary distribution of $S_t$ is $ \pi = (Z | Z > 0)$ for $Z \sim \mathrm{Poisson}(\lambda/\mu)$.
    %
    \[\pi_i = \frac{(\lambda/\mu)^i}{i!(e^{\lambda/\mu}-1)}, \quad \forall i \geq 1.\]
    %
\end{lemma}

\begin{proof}
    By detailed balance,
    %
    \[\pi_{i+1}(i+1)\mu=\lambda\pi_i,\qquad i\geq1.\]
    %
    \noindent Thus $\pi_i\propto(\lambda/\mu)^i/i!$, and normalization over $i\geq1$
    gives the stated zero-truncated Poisson distribution.
\end{proof}

\noindent The continuous-time Markov-chain ergodic theorem \cite[Section~20.3]{meyn2009_markov_chains_stochastic_stability} tells us as $T \to \infty$ the fraction of the time $\{S_t\}_{t \in [0,T]}$ spends in state 1
converges almost surely to $\pi_1 = \frac{\lambda / \mu}{e^{\lambda / \mu}-1}$. Heuristically, over long-branches our effective death rate then becomes
%
\[\mu_{\text{eff}} = \mu \left(1-\pi_1\right) = \mu \cdot \frac{e^{\lambda / \mu}-\lambda/\mu-1}{e^{\lambda / \mu}-1}.\]
%
\noindent In the low-polymorphism regime $\lambda \ll \mu$, then $\mu_{\text{eff}} \approx \frac{\lambda}{2} \ll \mu$. Thus
the original PCH model has a significantly smaller effective death rate than the coupled ID process.
If $\lambda \gg \mu$, then $\pi_1$ is exponentially small in $\lambda/\mu$, so the
process is rarely in state $1$, and it becomes well-approximated by our
ID process coupling, at least when homoplasy is also low.

Next, we show that our supports $|M_q^{(c)}|$ are reasonably concentrated about their means.
To this end, for $\alpha>0$ define the \emph{Orlicz pseudo-norm}
%
\[\|X\|_{\Psi_\alpha}:=\inf\left\{c>0:\E\exp\left[\left(\frac{|X|}{c}\right)^\alpha\right]\leq2\right\}.\]
%
\noindent A random variable with finite $\Psi_\alpha$ pseudo-norm is called
sub-Weibull$(\alpha)$. These pseudo-norms control tail decay and, through standard
concentration inequalities, deviations of sums of independent centered random variables
(e.g., \cite[Theorem~3.1]{kuchibhotla2022_beyond_sub_gaussian_high_dimensional_statistics}).

Here, we coarsely control the concentration of $|M_q|$ in terms of the number of states present at the root
and the relevant birth and death rates.

\begin{lemma}[Sub-Weibull Quartet Supports under PCH]
\label{lem:pch_subweibull_quartet_supports}
    Suppose $R:=\bigl\||c(r)|\bigr\|_{\Psi_1}<\infty$. Under the PCH model,
    %
    \[\bigl\||M_{uv||xy}^{(c)}|\bigr\|_{\Psi_{1/2}}
    \lesssim \left(1+R+\frac{\lambda}{\mu}\right)^2.\]
    %
    \noindent If the root cardinality has its stationary distribution, then the right-hand side is
    $\lesssim(1+\lambda/\mu)^2$.
\end{lemma}

\begin{proof}
    Along any root-to-leaf path $[r, \ell]$, couple the PCH cardinality process $S_t=|c(t)|$
    with $1+Z_t$, where $Z_t$ is an immigration-death process initialized at
    $Z_r=S_r-1$. At state $i$, the shifted process has death rate $(i-1)\mu$,
    whereas the PCH process has death rate $i\mu$ for $i>1$. A monotone coupling
    therefore gives $S_t\leq1+Z_t$.

    At any fixed $t$, $Z_t$ is the sum of a thinning of $S_r-1$ and an independent
    Poisson random variable with mean at most $\lambda/\mu$. The triangle inequality
    for $\Psi_1$, together with the Poisson bound
    $\|N\|_{\Psi_1}\lesssim1+\E N$, therefore gives
    %
    \[\bigl\||c(t)|\bigr\|_{\Psi_1}\lesssim
    1+R+\frac{\lambda}{\mu}.\]
    %
    Writing $U=c(u)$, and similarly for the other leaves, we have
    %
    \[|M_{uv||xy}^{(c)}|\leq |U\cap V|\,|X\cap Y|\leq |U|\,|X|.\]
    %
    \noindent Hence the H\"older-type inequality in
    \cite[equation~(3.5)]{kuchibhotla2022_beyond_sub_gaussian_high_dimensional_statistics}
    gives
    %
    \[\bigl\||M_{uv||xy}^{(c)}|\bigr\|_{\Psi_{1/2}}
    \leq \bigl\||U|\bigr\|_{\Psi_1}\bigl\||X|\bigr\|_{\Psi_1}
    \lesssim\left(1+R+\frac{\lambda}{\mu}\right)^2,\]
    %
    \noindent as desired. Under the stationary root law, $|c(r)|$ is zero-truncated
    $\mathrm{Poisson}(\lambda/\mu)$ and is stochastically dominated by
    $1+\mathrm{Poisson}(\lambda/\mu)$. Thus $R\lesssim1+\lambda/\mu$.
\end{proof}

 Later on, we will need to develop some control over the expected margin
 between the species tree topology and the other topologies. For
 $q \in A_s \setminus \{\cT_s^*\}$, define
    %
    \[D_q^{(c)} := |M^{(c)}_{\cT_s^*}| - |M_{q}^{(c)}|,
    \qquad \Delta_q^{(c)} := \E D_q^{(c)}.\]
    %
\noindent By Theorem~\ref{thm:correct_topo_max_expected_support},
$\Delta_q^{(c)}$ does not depend on the choice of incorrect topology $q$.
Let $\rho_N = \E |c(r) \cap \Sigma_N|$ be the
expected number of non-homoplastic states present at the root.
For a given (rooted) quartet, there are two possible topologies: a caterpillar tree and a balanced tree.
These are shown in Figure~\ref{fig:canby_quartet_topologies}.
Let $i,j,k$ denote the most recent common ancestors of these leaves, as depicted in the figure.

\begin{figure}[htbp]
    \centering
    \begin{subfigure}[t]{0.45\textwidth}
        \centering
        \begin{tikzpicture}[
          level distance=12mm,
          level 1/.style={sibling distance=24mm},
          level 2/.style={sibling distance=14mm},
          every node/.style={font=\small},
          edge from parent/.style={draw,-},
        ]
        \node (i) {$i$}
          child { node (j) {$j$}
            child { node (u) {$u$} }
            child { node (v) {$v$} }
          }
          child { node (k) {$k$}
            child { node (x) {$x$} }
            child { node (y) {$y$} }
          };
        \end{tikzpicture}
        \caption{Balanced quartet.}
    \end{subfigure}
    \hfill
    \begin{subfigure}[t]{0.45\textwidth}
        \centering
        \begin{tikzpicture}[
          level distance=12mm,
          level 1/.style={sibling distance=24mm},
          level 2/.style={sibling distance=14mm},
          every node/.style={font=\small},
          edge from parent/.style={draw,-},
        ]
        \node (i) {$i$}
          child { node (u) {$u$}}
          child { node (j) {$j$}
            child { node (v) {$v$} }
            child { node (k) {$k$}
                child { node (x) {$x$} }
                child { node (y) {$y$} }
            }
          };
        \end{tikzpicture}
        \caption{Caterpillar quartet.}
    \end{subfigure}
    \caption{The two rooted quartet topologies with split $uv||xy$. }
    \label{fig:canby_quartet_topologies}
\end{figure}


\begin{lemma}[Expected Margin Lower Bound]\label{lem:canby_expected_margin}

    If $\cT_s^*$ induces a balanced tree in $\cT^*$, then
    \[\Delta^{(c)}_q \geq \left(\frac{\lambda_N}{\mu} \right)^2(1-\eta_{ij})(1-\eta_{ik}) \eta_{uv} \eta_{xy}.   \]

\noindent If $\cT_s^*$ induces a caterpillar tree in $\cT^*$, then
%
\[\Delta^{(c)}_q \geq \left[\rho_N \eta_{ri} + \frac{\lambda_N}{\mu}(1-\eta_{ri})\right]
\frac{\lambda_N}{2\mu}(1-\eta_{jk})^2\eta_{uv}\eta_{xy}. \]
%

\end{lemma}

\noindent
The proof is given in the following subsection. These bounds highlight the quartets that are difficult to recover:
quartets where $u,v$ and $x,y$ are far from their respective most
recent common ancestors and where these common ancestors
are close together.


In both cases, the displayed margin lower bound is $\Omega(f^2)$ in the minimum branch length $f$,
assuming all other tree parameters, such as the tree height, are held fixed. To see why
this is natural, note that for a balanced quartet, the main outcome that produces support for the correct topology
but not the incorrect topologies is if $\alpha$ is born on $[i,j]$ and  $\beta$ is born on $[i,k]$. When these branches
are small, the probability each such birth occurs is $\Theta(f)$.

Likewise, for a caterpillar quartet, the main additional support for the correct topology comes from when
$\alpha$ dies along $[j,k]$ and $\beta$ is born along $[j,k]$. Again if branch $[j,k]$ is short, these events
both have probability of order $f$. Hence when $\cT^*$ has short branches, the expected margin for the true
topology can be quite small.


%%%%%%%%%%%%%%%%%%%%%%%%%%%%%%%%%%%%%%%%%%%%%%%%%%%%%%%%%%%%%%%%%%%%%%
% Proof of expected-margin bound
%%%%%%%%%%%%%%%%%%%%%%%%%%%%%%%%%%%%%%%%%%%%%%%%%%%%%%%%%%%%%%%%%%%%%%
\subsection{Proof of the Expected Margin Lower Bound}

\begin{proof}[Proof of Lemma \ref{lem:canby_expected_margin} (Expected support margin)]

Fix $q \in A_s \setminus \{\cT_s^*\}$ and recall
$D_q^{(c)}=|M_{\cT_s^*}^{(c)}|-|M_q^{(c)}|$ and $\Delta_q^{(c)}=\E D_q^{(c)}$.
For $a\in A_s$ and an ordered pair $(\alpha, \beta)$ of distinct non-homoplastic states, define
%
\[I_a(\alpha,\beta):=1\{(\alpha,\beta)\in M_a^{(c)}\}.\]
%
\noindent\emph{Balanced case.}
Suppose first that $\cT_s^*=uv||xy$ induces a balanced tree in $\cT^*$, with internal
vertices $i,j,k$ as in Figure~\ref{fig:canby_quartet_topologies}.
Let $I := c(i)$ be the set of micro-states at vertex $i$, and likewise for $J,K$.
Expanding $D_q^{(c)}$ as a sum over ordered pairs of states, the cancellation argument in the proof of
Theorem 1 from~\cite{evans2026_canby_identifiability} pairs contributions that support an
incorrect split and leaves the following nonnegative contributions:
%
\[\E [D_q^{(c)} \mid J, K] = \sum_{(\alpha, \beta) \in \cA(J, K)} \E(I_{uv||xy}(\alpha, \beta) \mid J, K),\]
%
\noindent where
%
\[\cA(J, K) = \{(\alpha, \beta) : \alpha\neq\beta \in \Sigma_N,\ \alpha \in J,\ \beta \in K,\ \{\alpha, \beta\} \nsubseteq J \cap K\}.\]
%
\noindent We see that
%
\[\cB(J, K) := \{(\alpha, \beta) : \alpha \in J \setminus K,\ \beta \in K \setminus J\} \subseteq \cA(J, K).\]
%
\noindent For $(\alpha, \beta) \in \cB(J, K)$, we just need $\alpha, \beta$ to survive until the leaves, so
%
\[\E(I_{uv||xy}(\alpha, \beta) \mid J, K) \geq \eta_{ju} \eta_{jv} \eta_{kx} \eta_{ky}.\]
%
\noindent Hence
%
\[\Delta_q^{(c)} \geq \E |\cB(J, K)| \cdot \eta_{ju} \eta_{jv} \eta_{kx} \eta_{ky}.\]
%
\noindent We lower-bound $|\cB(J,K)|$ using only pairs for which state $\alpha$ is born
along $(i,j]$ and survives until $j$, while state $\beta$ is born along $(i,k]$ and
survives until $k$. In this case the roles of $J,K$ decouple. By the coupling and
Poisson thinning, the numbers of such $\alpha$ and $\beta$ stochastically dominate independent
$\mathrm{Poisson}(m_0(\ell_{ij}))$ and $\mathrm{Poisson}(m_0(\ell_{ik}))$ random variables,
respectively, where $m_0(t) = \frac{\lambda_N}{\mu}(1-\eta_t)$. This produces on average
$m_0(\ell_{ij}) \cdot m_0(\ell_{ik})$ pairs $(\alpha, \beta)$.
Thus
%
\[\E|\cB(J,K)|\geq m_0(\ell_{ij})m_0(\ell_{ik}),\]
%
\noindent which gives the desired bound on $\Delta_q^{(c)}$ when the induced subtree is balanced.

\noindent\emph{Caterpillar case.}
Suppose instead that $\cT_s^*=uv||xy$ induces the caterpillar tree in
Figure~\ref{fig:canby_quartet_topologies}. The analogous cancellation argument from the proof of
Theorem 1 in~\cite{evans2026_canby_identifiability} again cancels the contributions
supporting an incorrect split and gives
%
\[\E [D_q^{(c)} \mid I, J, K] = \sum_{(\alpha, \beta) \in \tilde{\cA}(I, J, K)} \E(I_{uv||xy}(\alpha, \beta) \mid I, J, K),\]
%
\noindent where
\[\tilde{\cA}(I, J, K) = \{(\alpha, \beta) : \alpha\neq\beta \in \Sigma_N,\ \alpha \in I \cap J,\ \beta \in K,\ \{\alpha, \beta\} \nsubseteq J \cap K\}.\]
%
\noindent Again, we see this set contains
%
\[\tilde{\cB}(I, J, K) := \{(\alpha, \beta) : \alpha \in (I \cap J) \setminus K,\ \beta \in K \setminus J\}.\]
%
For $(\alpha, \beta) \in \tilde{\cB}(I, J, K)$ we see $(\alpha,\beta) \in M_{uv||xy}$ if and only if:

\begin{itemize}
    \item $\alpha$ survives on $(i, u]$ and $(j, v]$.
    \item $\beta$ survives on $(k,x]$ and $(k,y]$.
\end{itemize}

\noindent so the survival bound above gives
%
\[\E(I_{uv||xy}(\alpha, \beta) \mid I, J, K) \geq \eta_{iu}\eta_{jv} \eta_{kx} \eta_{ky}.\]
%
\noindent In particular, $(\alpha, \beta) \in \tilde{\cB}(I,J,K)$ can occur if
%
\begin{itemize}
    \item $\alpha$ is born before $i$, survives along $(i,j]$, and dies on $(j,k]$.
    \item $\beta$ is born on $(j,k]$ and survives until $k$.
\end{itemize}

\noindent Let $\rho_N := \E |c(r) \cap \Sigma_N|$ be the expected number of
non-homoplastic states present at the root. Clearly we have
%
    \[\E |I\cap\Sigma_N| \geq \rho_N \eta_{ri} + m_0(\ell_{ri})\]
%
\noindent as non-homoplastic states in $I$ can come either from the root or be born along $[r,i]$. Of these, at
least
%
    \[\E |I \cap J\cap\Sigma_N| \geq \left(\rho_N \eta_{ri} + m_0(\ell_{ri})\right) \eta_{ij}. \]
%
\noindent survive until they reach $j$, giving our candidates for $\alpha$. Condition on $I \cap J\cap\Sigma_N$ and
fix $\alpha \in I \cap J\cap\Sigma_N$. Let
%
    \[X_\alpha := \# \{\beta : (\alpha, \beta) \in \tilde{\cB}(I, J, K)\}\]
%
\noindent be the number of pairs that $\alpha$ contributes to $\tilde{\cB}$. If a non-homoplastic state
$\beta$ is born at location $s \in [j, k]$, it appears in $X_\alpha$ with probability at least
%
    \[p(s):= \exp(-\mu \ell_{sk} ) \cdot (1- \exp(-\mu \ell_{sk})),\]
%
\noindent Indeed, $\beta$ survives $[s,k]$ with probability $e^{-\mu\ell_{sk}}$.
If $\alpha$ is still present at $s$, the presence of $\beta$ ensures that a ring of $\alpha$'s
death clock in $[s,k]$ removes it; if $\alpha$ has already died, the desired event already holds.
If $N$ represents the coupled rate-$\lambda_N$ Poisson process of
non-homoplastic births in $[j,k]$, which is independent of the death clocks, then Campbell's formula gives
%
    \[\E X_\alpha \geq \E \left(\sum_{s \in N} p(s)\right) = \lambda_N \int_{0}^{\ell_{jk}} \exp(-\mu s) (1 - \exp(-\mu s)) \; ds = \frac{\lambda_N}{2\mu} (1-\exp(-\mu \ell_{jk}))^2. \]
%
\noindent Hence we see
%
    \[\E |\tilde{\cB}(I,J,K)| \geq \left[\rho_N \eta_{ri} + \frac{\lambda_N}{\mu}(1-e^{-\mu \ell_{ri}})\right] \eta_{ij} \cdot \left[\frac{\lambda_N}{2\mu} (1-\eta_{jk})^2\right],\]
%
\noindent giving the desired bound.
\end{proof}

%%%%%%%%%%%%%%%%%%%%%
% properties of ID Process
%%%%%%%%%%%%%%%%%%%%%
\subsection{Properties of $ID(\theta, \lambda, \mu; \cT^*)$}

We first record several elementary properties of these immigration-death processes that
will be used throughout the chapter. Along a single branch $[r,\ell]$, the cardinality
process is always Poissonian with an explicit mean.

\begin{lemma}[Population Sizes in ID Processes]\label{lem:pop_sizes_in_ID_processes}
    If $(Y_t)_{t \geq 0}$ is the cardinality process along a single branch $[r, \ell]$ of $ID(\theta, \lambda, \mu; \cT^*)$, then
    %
    \[|Y_t| \sim \mathrm{Poisson}(m_\theta(t)), \quad m_\theta(t) = \theta e^{-\mu t} + \frac{\lambda}{\mu}(1-e^{-\mu t}).\]
    %
    \noindent In particular, its stationary distribution is $\mathrm{Poisson}(\lambda/\mu)$.
\end{lemma}

\begin{proof}
    The Poisson probabilities $p_k(t)=e^{-m_\theta(t)}m_\theta(t)^k/k!$ satisfy the Kolmogorov forward equations
    %
    \[\frac{d}{dt}p_k = \lambda p_{k-1} + \mu (k+1) p_{k+1} - (k\mu + \lambda) p_k,
    \qquad p_k(0)=\frac{e^{-\theta}\theta^k}{k!}.\]
    %
    Hence $|Y_t|\sim\mathrm{Poisson}(m_\theta(t))$.
\end{proof}

Next, the quartet supports are products of independent Poisson random variables.

\begin{lemma}[Quartet Supports under ID Process]
\label{lem:quartet_supports_under_ID_process}
    Suppose $Y_t \sim ID(\theta, \lambda, \mu; \cT^*)$. For any subset $\mathcal U$
    of the leaves of $\cT^*$ and each $A\subseteq\mathcal U$, define
    %
    \[N_A^{\mathcal U}:=\left|\left\{\alpha:
    \alpha\in Y_\ell\ \text{for every }\ell\in A,
    \ \alpha\notin Y_\ell\ \text{for every }\ell\in\mathcal U\setminus A
    \right\}\right|.\]
    %
    \noindent Then $\{N_A^{\mathcal U}\}_{A\subseteq\mathcal U}$ are mutually independent Poisson random variables.
\end{lemma}

\begin{proof}
    Regard the states born on $\cT^*$ as a Poisson point process with intensity
    $\theta\delta_r(dz)+\lambda\,dz$, and mark each state by the subset of $\mathcal U$ on which it appears.
    Kingman's marking theorem implies that the counts associated with distinct marks are independent Poisson random variables
    \cite{kingman1992_poisson_processes}.
\end{proof}

Choosing $\cU=\{u,v,x,y\}$ gives the following immediate corollary.

\begin{corollary}\label{cor:id_process_supports_poisson_products}
    Under the standalone $ID(\theta, \lambda, \mu; \cT^*)$ process, the quartet supports equal
    %
    \[|M_{uv||xy}| = N_{uv}N_{xy}, \quad |M_{ux||vy}| = N_{ux}N_{vy}, \quad |M_{uy||vx}| = N_{uy}N_{vx},\]
    %
    \noindent for $\{N_{uv},\ldots,N_{vx}\}$ mutually independent Poisson random variables.
\end{corollary}

\noindent The product-Poisson representation sharpens
Lemma~\ref{lem:pch_subweibull_quartet_supports} under the ID approximation by retaining
the pair-specific survival factors.

\begin{lemma}[Sub-Weibull Quartet Supports]
\label{lem:canby_subweibull_quartet_supports}
    Let $\zeta = \theta \vee \frac{\lambda}{\mu}$. Under $ID(\theta, \lambda, \mu; \cT^*)$,
    $|M_{uv||xy}|$ is sub-Weibull$(1/2)$ and
    %
    \[\bigl\||M_{uv||xy}|\bigr\|_{\Psi_{1/2}} \lesssim
    \left(1 + \zeta\eta_{uv} \right) \left(1 + \zeta\eta_{xy} \right).\]
    %
\end{lemma}

\begin{proof}
    By Corollary~\ref{cor:id_process_supports_poisson_products},
    $|M_{uv||xy}|=N_{uv}N_{xy}$ for independent Poisson random variables $N_{uv}$ and $N_{xy}$.
    Lemma~\ref{lem:pop_sizes_in_ID_processes} and Poisson thinning give
    $\E N_{uv}\leq\zeta\eta_{uv}$ and $\E N_{xy}\leq\zeta\eta_{xy}$.
    Hence the H\"older-type inequality in~\cite[equation~(3.5)]{kuchibhotla2022_beyond_sub_gaussian_high_dimensional_statistics} gives
    %
    \[\bigl\||M_{uv||xy}|\bigr\|_{\Psi_{1/2}} \leq \|N_{uv}\|_{\Psi_1}\|N_{xy}\|_{\Psi_1}.\]
    %
    \noindent For $N\sim\mathrm{Poisson}(m)$ with $m>0$,
    %
    \[\|N\|_{\Psi_1}=\frac{1}{\log(1+\log(2)/m)}\leq 1+\frac{m}{\log(2)},\]
    %
    \noindent while the case $m=0$ is trivial. The result follows.
\end{proof}

Beyond developing concentration inequalities, Corollary~\ref{cor:id_process_supports_poisson_products}
allows us to bound the expected margin
%
    \[\Delta_q := \E [ |M_{\cT_s^*}| - |M_q|],
    \qquad q \in A_s \setminus \{\cT_s^*\},\]
%
\noindent between the correct topology $\cT_s^*$ and any incorrect topology $q \in A_s \setminus \{\cT_s^*\}$.
If, for example, $\cT_s^* = uv || xy$ and $q = ux||vy$, then writing $m_{ab}:=\E N_{ab}$ we see
%
\[\Delta_q =m_{uv}m_{xy}-m_{ux}m_{vy}.\]
%
Lemma~\ref{lem:pop_sizes_in_ID_processes} then gives the simple upper bound
%
\[m_{ab}\leq\left(\theta\vee\frac{\lambda}{\mu}\right)\eta_{ab},\]
%
\noindent which can be used to upper bound the margin $\Delta_q$. The corresponding
expected-margin lower bounds are as follows.

\begin{corollary}[Expected Margin under the ID Process]
\label{cor:expected_margin_under_id}
Suppose characters are generated under $ID(\theta,\lambda,\mu;\cT^*)$. If $\cT_s^*$ induces a balanced quartet, then
\[
    \Delta_q \geq \left(\frac{\lambda}{\mu}\right)^2
    (1-\eta_{ij})(1-\eta_{ik})\eta_{uv}\eta_{xy}.
\]
If $\cT_s^*$ induces a caterpillar quartet, then
\[
    \Delta_q \geq
    \left[\theta\eta_{ri}+\frac{\lambda}{\mu}(1-\eta_{ri})\right]
    \frac{\lambda}{\mu}(1-\eta_{jk})^2\eta_{uv}\eta_{xy}.
\]
\end{corollary}

\begin{proof}
The balanced bound follows from the proof of Lemma~\ref{lem:canby_expected_margin} with
$\lambda_N=\lambda$ and $\rho_N=\theta$. For a caterpillar quartet, a candidate state $\alpha$
dies along $[j,k]$ with probability $1-\eta_{jk}$, independently of the number of states $\beta$
born along this branch and surviving to $k$, whose mean is
$m_0(\ell_{jk})=\lambda\mu^{-1}(1-\eta_{jk})$. Multiplying these quantities and the relevant
survival probabilities gives the result.
\end{proof}

\section{A Restricted Search-Space}
\label{sec:bipartition_covers_canby}

\noindent The exact version of ASTRAL~\eqref{eqn:astral_quartet_support_optimization_problem}
becomes intractable to solve even for trees with a modest number of leaves. Instead, we must construct a
set $\cX$ of bipartitions which will limit the search space for the true tree $\cT^*$, which is the content of this section.

To this effect, one useful class of sets when it comes to building $\cX$ are the leaf sets
%
\[L_\alpha := \{\ell \in \cL(\cT^*) : \alpha \in c(\ell)\}.\]
%
\noindent Say a subset $L \subseteq \cL(\cT^*)$ of leaves is \textit{contained} within a given bipartition $\phi := (A|B)$ if
%
\[L \subseteq A \qquad \text{or} \qquad L \subseteq B.\]
%
 If $\alpha \in \Sigma_N$ is born
 on an edge $e$ inducing bipartition $\phi_e = (A_e | B_e)$, then $L_\alpha$ is contained in $\phi_e$. Moreover, if no
deaths of $\alpha$ occur after it is born then in fact $L_\alpha \in \{A_e, B_e\}$. Hence
one choice for the candidate bipartitions is
%
\[\cX_1 := \left\{\left(L^{(c)}_\alpha \mid \cL(\cT^*) \setminus L^{(c)}_\alpha\right) :
c \in \cC,\ \alpha \in \Sigma^{(c)},\ \varnothing \subsetneq L^{(c)}_\alpha \subsetneq \cL(\cT^*)\right\}.\]
%
\noindent At stationarity, we typically have around $\lambda \mu^{-1}$ states in each leaf,
so roughly we can see $|\cX_1| \lesssim k |\cC| \lambda \mu^{-1}$ grows only linearly in $k$.
Under mild assumptions this alone is sufficient for ensuring statistical consistency,
since with positive probability a given non-homoplastic $\alpha$ is born
along an edge $e$ and does not die at all. More precisely,

% bipartition cover bound
\begin{theorem}[Bipartition Cover Bounds for PCH With No Borrowing]
\label{thm:canby_bipartition_cover_bounds_no_borrowing}
Let $\cB$ be the set of nontrivial bipartitions of a phylogenetic tree $\cT^*$ with $k$ leaves
and minimum/maximum branch lengths $f_{\min}, f_{\max}$, respectively. Suppose further
we have independent PCH characters with parameters
$\left\{\Sigma^{(c)}, \lambda^{(c)}, \mu^{(c)}, \nu^{(c)}, \rho^{(c)}\right\}_{c \in \cC}$
on $\cT^*$ with no borrowing and whose cardinality processes have immigration-death rates. Let
%
\[\cC_N := \{c \in \cC \; : \; \nu^{(c)}(\Sigma_N^{(c)}) > 0\},\]
%
\noindent be the set of characters with positive chance of non-homoplasy.
 Furthermore, assume we have the following assumptions on the parameters of the PCH model:

\begin{enumerate}

    \item[(i)] \textit{Nonvanishing Birth Rates:} Suppose that $\lambda^{(c)} \nu^{(c)}(\Sigma_N^{(c)}) \geq b > 0$ for all $c \in \cC_N$.

    \item[(ii)] \textit{Bounded death rates}: Suppose $\mu^{(c)} \leq D < \infty$ for all $c \in \cC_N$.
\end{enumerate}
%
\noindent Then the coverage failure probability decays exponentially:
%
\[\P(\cB \nsubseteq \cX_1)  \leq\sum_{\ell=2}^{k-2} q_\ell^{|\cC_N|} \leq (k-3)q_{k-2}^{|\cC_N|},\qquad q_\ell := \exp\left(-bf_{\min} \exp\left[-2\ell Df_{\max}\right]\right).\]
%
\noindent In particular, we have asymptotic consistency when $|\cC_N| \to \infty$
%
\[\lim_{|\cC_N| \to \infty} \P(\cB \subseteq \cX_1) = 1.\]
%
\end{theorem}
%
\begin{proof}
Fix an edge $e \in E(\cT^*)$ inducing a nontrivial bipartition, with $k_e$ descendants, and a given character $c$.
For ease of notation, we drop all the $(c)$ superscripts for now. If $\phi_e$ is the
corresponding bipartition, we will lower bound the probability that
$\phi_e = (L_\alpha \mid \cL(\cT^*) \setminus L_\alpha)$ for some $\alpha \in \Sigma_N$.

First, note by our coupling to the ID process, the number $N_e$ of non-homoplastic
states born along the edge $e$ satisfies
%
\[N_e \succeq_{st} \mathrm{Poisson}(\lambda T_e \cdot \nu(\Sigma_N)) \succeq_{st} \mathrm{Poisson}(\eta), \qquad \eta := \lambda f_{\min} \nu(\Sigma_N),\]
%
\noindent where $T_e$ is the length of the branch $e$ and $f_{\min}$ is the minimum branch length in the tree.

For a given non-homoplastic $\alpha$ born along the edge $e$, we can lower bound the probability it has no deaths below $e$. The probability $\alpha$ survives along a branch of length $T$ is at least $\exp(-\mu T)$. There are $2k_e - 1 \leq 2k_e$ edges (including $e$) that $\alpha$ can die on. Each is at most length $f_{\max}$, so we can lower bound the probability $\alpha$ has no deaths at all by
%
\[\P(\alpha \text{ has no deaths} \mid \alpha \in \Sigma_N,\ \alpha \text{ born on } e) \geq \exp(- \tau_{k_e}), \qquad \tau_{k_e} := 2\mu k_e f_{\max}.\]
%
\noindent Using this result followed by $N_e \succeq_{st} \mathrm{Poisson}(\eta)$ we see
%
\[\P\bigl(\exists\alpha \; \phi_e = (L_\alpha \mid \cL(\cT^*) \setminus L_\alpha)\bigr) \geq \E_{n \sim N_e} \left(1-(1-\exp(-\tau_{k_e}))^n\right) \geq 1-\exp(-\eta e^{-\tau_{k_e}}).\]
%
From here, the caterpillar bound from~\cite{mcnulty2026_improved_bipartition_cover_bound} applied to our $|\cC|$ characters implies
%
\[\P(\cB \subseteq \cX_1) \geq 1 - \sum_{\ell=2}^{k-2} \prod_{c \in \cC}\exp\left(-\eta^{(c)} e^{-\tau_{\ell}^{(c)}}\right) =  1 - \sum_{\ell=2}^{k-2} \exp\left(-\sum_{c \in \cC}\eta^{(c)} e^{-\tau_{\ell}^{(c)}}\right).\]
%
\noindent Note none of the above uses our boundedness assumptions on the birth or death rates $\lambda^{(c)}, \mu^{(c)}$. Using these assumptions on the birth and death rates, we see
%
\[\P(\cB \subseteq \cX_1) \geq 1 - \sum_{\ell=2}^{k-2}q_\ell^{|\cC_N|} \geq 1-(k-3)q_{k-2}^{|\cC_N|}, \qquad q_\ell := \exp\left(-bf_{\min} \exp\left[-2\ell Df_{\max}\right]\right).\]
%
\noindent Hence the coverage failure probability decays exponentially.
\end{proof}

In particular, running ASTRAL as proposed in \cite{evans2026_canby_identifiability},
with its search restricted to $\cX_1$, gives a statistically consistent estimator of
$\cT^*$. The bound requires $\exp(O(k))$ characters to achieve a bipartition cover
with high probability. This is essentially the same rate
bipartition-covers achieve under the standard multispecies-coalescent model to which
ASTRAL is typically applied \cite{uricchio2016_bipartition_cover}, \cite{mcnulty2026_improved_bipartition_cover_bound}.

This result should be viewed as a baseline consistency guarantee rather
than a practical sample-complexity result. The set $\cX_1$ is poor at generating
large bipartitions: if $\phi_e=(A|B)$, the probability that a state born on edge $e$
has no subsequent deaths is exponentially small in $|A|$. We therefore propose two
heuristics for expanding $\cX_1$. One option is to allow multiple states to contribute:
%
\[\cX_m := \left\{(A \mid \cL(\cT^*) \setminus A) \; : \;
A = \bigcup_{i=1}^r L_{\alpha_i}^{(c)},\ c \in \cC,\ \alpha_i \in \Sigma^{(c)},\
1 \leq r \leq m,\ \varnothing \subsetneq A \subsetneq \cL(\cT^*)\right\}.\]
%
\noindent It is easy to see that $\cX_{m_1} \subseteq \cX_{m_2}$ for $m_1 \leq m_2$, so in
 particular $\cX_1 \subseteq \cX_m$. Hence asymptotic consistency using $\cX_m$ follows
 from that of $\cX_1$. To avoid combining leaf sets from
 completely different regions in the tree, we can also impose a compatibility condition such as
 $L_{\alpha_i}^{(c)} \cap L_{\alpha_j}^{(c)} \neq \emptyset$.

Alternatively, for a nonnegative integer $r$, we can expand each observed leaf set by adding at most $r$ leaves. Define
%
\[\cX_{\leq r} := \left\{(A \mid \cL(\cT^*) \setminus A) \; : \;
\exists\, c \in \cC,\ \alpha \in \Sigma^{(c)}\text{ such that }
L_\alpha^{(c)} \subseteq A \subsetneq \cL(\cT^*),\ |A \setminus L_\alpha^{(c)}| \leq r\right\}.\]
%
\noindent For a fixed leaf set $L_\alpha^{(c)}$, this generates at most
$\sum_{j=0}^r \binom{k-|L_\alpha^{(c)}|}{j} = O(k^r)$ candidate bipartitions when $r$ is fixed.
Moreover, $\cX_1 \subseteq \cX_{\leq r}$, so again this expansion preserves the asymptotic-consistency guarantee above.

The two expansions address different failure modes. The set $\cX_m$ combines observed
state patterns and is therefore more closely tied to the evolutionary model, but if character
$c$ produces $h_c$ distinct leaf sets, it may generate $O(h_c^m)$ unions; compatibility
constraints can reduce this growth. In contrast, $\cX_{\leq r}$ has the predictable
$O(k^r)$ cost per observed leaf set, but adds leaves without using character-level evidence.
Thus $\cX_m$ is more model-informed but potentially more combinatorial, whereas
$\cX_{\leq r}$ offers a simpler computational budget. We view both as heuristics, since
we do not presently provide finite-sample coverage bounds for either expansion.

%%%%%%%%%%%%%%%%%%%%%%%%%%%%%%%%%%%
% PCH-ASTRAL Sample Complexity
%%%%%%%%%%%%%%%%%%%%%%%%%%%%%%%%%%%
\section{PCH-ASTRAL Sample-Complexity Bounds}
\label{sec:pch_astral_sample_complexity}

In this section, we bound the number of independent characters required by PCH-ASTRAL
to recover the true species tree $\cT^*$ with high probability. The result applies to both
the original PCH model and its coupled $ID(\theta,\lambda,\mu;\cT^*)$ approximation; it only
requires uniform sub-Weibull control of the quartet supports.

\begin{lemma}[Sample Complexity PCH-ASTRAL]
\label{lem:sample_complexity_pch_astral}
    Suppose $\{c_i\}_{i=1}^n$ are iid characters generated under either the PCH model or
    $ID(\theta,\lambda,\mu;\cT^*)$. Define
    %
    \[\kappa:=\max_{s\in{[k]\choose4}}\max_{q\in A_s}
    \bigl\||M_q|\bigr\|_{\Psi_{1/2}}<\infty.\]
    %
    \noindent Define our worst-case margin
        \[
            L:=\min_{s\in {[k]\choose 4}}\min_{q\in A_s\setminus\{\cT_s^*\}}\E\left[|M_{\cT_s^*}|-|M_q|\right]>0.
        \]
    \noindent Then if we have at least
    %
    \[n \gtrsim \max\left\{\frac{\kappa^2}{L^2} t_\delta, \; \frac{\kappa}{L} t_\delta^2 \right\}, \quad t_\delta := \log\left(\frac{4{k \choose 4}}{\delta} \right), \]
    %
    \noindent characters, we recover the true species tree with probability at least $1-\delta$.
\end{lemma}

\begin{proof}[Proof of Lemma \ref{lem:sample_complexity_pch_astral} (Sample Complexity PCH-ASTRAL)]
    If $\hat{q}_s = \cT_s^*$ for every leaf-set $s$, then PCH-ASTRAL will trivially return the correct species tree
    $\cT^*$, as the set of topologies $\{\hat{q}_s\}$ will exactly match those generated by the species tree $\cT^*$.
    Hence by union bound
    %
     \[\P(\hat{\cT} = \cT^*) \geq 1 - \sum_{s} \P(\hat{q}_s \neq \cT_s^*).\]
    %
    \noindent Hence it suffices to upper bound $p_s := \P(\hat{q}_s \neq \cT_s^*)$. Note
    %
    \[p_s \leq \sum_{q \in A_s \setminus \{\cT_s^*\}} \P(S_{\cT_s^*} \leq S_q).\]
    %
    \noindent so it suffices to bound each term on the right hand side. For $q \in A_s \setminus \{\cT_s^*\}$, define
    %
        \[D_i := |M^{(c_i)}_{\cT_s^*}| - |M^{(c_i)}_q|, \qquad S_{\cT_s^*} - S_q = \frac{1}{n} \sum_{i=1}^n D_i.\]
    %
    \noindent Since $L \leq \E D_i$ we have
    %
        \[\P(S_{\cT_s^*} \leq S_q) \leq \P\left(\Bigg| \sum_{i=1}^n (D_i - \E D_i)\Bigg| \geq n L\right).\]
    %
    \noindent The definition of $\kappa$, together with the centering and quasi-triangle
    inequalities for $\Psi_{1/2}$, gives
    %
        \[\|D_i - \E D_i\|_{\Psi_{1/2}} \lesssim
        \bigl\||M_{\cT_s^*}|\bigr\|_{\Psi_{1/2}}
        + \bigl\||M_q|\bigr\|_{\Psi_{1/2}} \lesssim \kappa.\]
    %
    \noindent Applying the concentration inequality~\cite[Theorem 3.1]{kuchibhotla2022_beyond_sub_gaussian_high_dimensional_statistics} gives
    %
        \[\P\left(\Bigg| \sum_{i=1}^n (D_i - \E D_i)\Bigg| \gtrsim \kappa( \sqrt{nt} + t^2)\right) \leq 2e^{-t}.\]
    %
    \noindent Choose $t_\delta = \log(4{k \choose 4} / \delta)$. Note our assumptions on $n$ imply
    %
    \[nL \gtrsim \kappa (\sqrt{nt_\delta} + t_\delta^2),\]
    %
    \noindent and thus we have
    %
        \[\P(S_{\cT_s^*} \leq S_q) \leq 2\exp(-t_\delta).\]
    %
    \noindent Taking a union bound over quartets and the two incorrect topologies then gives
    %
        \[\P(\hat{\cT} = \cT^*) \geq 1 - 2{k \choose 4} \cdot 2\exp(-t_\delta) = 1-\delta,\]
    %
    \noindent as desired.
\end{proof}

\noindent Thus $\kappa$ is the worst-case Orlicz scale of the quartet supports.
For the ID approximation, Lemma~\ref{lem:canby_subweibull_quartet_supports} gives, for
$\zeta=\theta\vee\lambda/\mu$,
%
\[\kappa\lesssim\max_{s\in{[k]\choose4}}\max_{uv||xy\in A_s}
(1+\zeta\eta_{uv})(1+\zeta\eta_{xy})\leq(1+\zeta)^2.\]
%
\noindent Under the PCH model, Lemma~\ref{lem:pch_subweibull_quartet_supports} gives
%
\[\kappa\lesssim\left(1+\bigl\||c(r)|\bigr\|_{\Psi_1}+\frac{\lambda}{\mu}\right)^2.\]
%
\noindent Under the stationary root law, this reduces to
$\kappa\lesssim(1+\lambda/\mu)^2$.
%
\noindent The worst-case margin $L$ is the main remaining parameter of interest.
Corollary~\ref{cor:expected_margin_under_id} lower bounds $L$ under $ID(\theta, \lambda, \mu; \cT^*)$
purely in terms of the underlying model parameters. While this bound inherently depends on the topology of $\cT^*$,
through some coarse analysis we can extend it to a topology-free bound. Suppose
%
    \[d := \max_{uv||xy \in \cQ(\cT^*)} \ell_{uv} \vee \ell_{xy}\]
%
\noindent is the maximal distance between any two leaves on the same side of a quartet.
Then the margin lower bounds in Corollary~\ref{cor:expected_margin_under_id} roughly become
%
    \[\Delta_q \geq \left(\frac{\lambda}{\mu}\right)^2 (1-e^{-\mu f})^2 (e^{-\mu d})^2 \gtrsim \lambda^2 f^2 e^{-2\mu d}, \]
%
\noindent when the induced subtree in $\cT^*$ is balanced and
%
        \[\Delta_q \gtrsim \lambda\mu f^2 e^{-2\mu d} \left(\theta \wedge \frac{\lambda}{\mu}\right), \]
%
\noindent when the induced subtree is a caterpillar. Thus, for a tree of fixed height and fixed model parameters,
$\Delta_q = \Omega(f^2)$, and so our sample-complexity bound becomes
%
    \[n \gtrsim \max\left\{f^{-4} \log(k), \; f^{-2} \log^2(k)\right\}.\]
%
\noindent This remains significantly worse than the $\Theta(f^{-2} \log(k))$ of ASTRAL under the multi-species coalescent model~\cite{shekhar2018_how_many_genes_are_enough}. This bound is likely loose, and a sharper analysis may improve the rate.



%%%%%%%%%%%%%%%%%%%%%%%%%%%%%%%%%%%%%%%%%%%%%%%%
% Alternative ASTRAL procedure
%%%%%%%%%%%%%%%%%%%%%%%%%%%%%%%%%%%%%%%%%%%%%%%%
\section{Alternative ASTRAL Procedure}
\label{sec:canby_alternative_astral_procedure}

Another inference method, closer in spirit to the original ASTRAL framework, is to let each
character independently vote on the topology of each quartet. These character-level quartet
votes then play the role of the quartet information induced by individual gene trees in the
traditional multispecies coalescent setting.

More formally, for each character $c_i$ and each leaf set
$s=\{u,v,x,y\}$, define
\[
    \tau_s^{(c_i)}
    :=
    \arg\max_{q\in A_s} |M_q^{(c_i)}|, \quad I^{(c_i)}_s := \left\{\max_{q\in A_s} |M_q^{(c_i)}| > 0\right\},
\]
\noindent with ties broken uniformly. We then solve~\eqref{eqn:astral_quartet_support_optimization_problem} with
quartet supports
%
    \[W(q) = \sum_{i=1}^n 1\{\tau_s^{(c_i)} = q\} \cdot 1\{I_s^{(c_i)}\}.\]
%
The indicator $1\{I_s^{(c_i)}\}$ essentially just tells us to ignore characters which provide no information
about the given quartet. We call this single-character voting procedure \textit{PCH-ASTRAL(sc)}.
This inference scheme is only really reasonable if for
each quartet, the correct topology is the most likely one. Namely

\begin{equation}
    \P(\tau_s^{(c)} = \cT_s^*) > \P(\tau_s^{(c)} = q) \qquad q \neq \cT_s^*.
    \label{eqn:anomaly_free_property}
\end{equation}

\noindent On $I_s^c$, all three supports vanish, so uniform tie-breaking assigns each
topology probability $1/3$. Consequently, property~\eqref{eqn:anomaly_free_property}
is equivalent to the corresponding strict inequalities conditional on $I_s$.

If property~\eqref{eqn:anomaly_free_property} holds for every quartet $\cT_s^* \in \cQ(\cT^*)$,
we say the PCH model is in the \textit{anomaly-free zone}. When this occurs, the above
ASTRAL-based procedure is asymptotically consistent and has sample-complexity essentially 
quadratic in the margin. Note the following makes no assumptions on the distribution of the characters
beyond the anomaly-free property, so it applies equally to the PCH model and the standalone $ID(\theta, \lambda, \mu; \cT^*)$ process.

\begin{theorem}[Sample Complexity PCH-ASTRAL(sc)]
\label{thm:sample_complexity_separate_characters}
    Suppose we have a collection of $n$ iid characters $\{c_i\}$ in the anomaly-free zone, and let $0<\varepsilon<1$. If
        \[n > \frac{8}{p_{\min}\Delta^2} \log\left(\frac{3\binom{k}{4}}{\varepsilon}\right),\]
        %
        \noindent where
        %
        \[\Delta := \min_{s \in {[k] \choose 4}} \min_{q \in A_s\setminus\{\cT_s^*\}} \left[\P(\tau_s^{(c)} = \cT_s^* \mid I_s) - \P(\tau_s^{(c)} = q \mid I_s) \right], \quad p_{\min} := \min_{s \in {[k] \choose 4}} \P(I_s),\]
        %
    \noindent then PCH-ASTRAL(sc) recovers $\cT^*$ with probability at least $1-\varepsilon$.
\end{theorem}

Note $\Delta$ gives us some notion of the margin the correct topology has over the incorrect ones. On the other hand, $p_{\min}$
measures how many informative characters we tend to observe for our worst-case quartet, so it gives some notion of
the effective sample size. The argument is analogous to the proof of Theorem 2.1 in~\cite{shekhar2018_how_many_genes_are_enough}.

\begin{proof}
Let $K:=\binom{k}{4}$. For a fixed leaf set $s$, write
$p_s:=\P(I_s)$ and let $n_s$ be the number of informative characters. Since
$n_s\sim\mathrm{Binomial}(n,p_s)$, Chernoff's bound gives
\[
    \P\left(n_s<\frac{np_s}{2}\right)
    \leq \exp\left(-\frac{np_{\min}}{8}\right).
\]
Therefore, by union bound
\[
    \P\left(\min_{s\in {[k]\choose 4}}n_s<\frac{np_{\min}}{2}\right)
    \leq K\exp\left(-\frac{np_{\min}}{8}\right).
\]

\noindent Hence with high probability, every leaf set $s$ has at least
$np_{\min}/2$ informative characters.

Now fix $s$ and label its two incorrect topologies $q_1,q_2$. Conditional on
$I_s$, define
\[
    p_s^*:=\P(\tau_s^{(c)}=\cT_s^*\mid I_s),
    \qquad
    r_{s,j}:=\P(\tau_s^{(c)}=q_j\mid I_s),\quad j=1,2.
\]
For each informative character $c_i$, let
\[
    Z_i^{(j)}
    :=1\{\tau_s^{(c_i)}=\cT_s^*\}
    -1\{\tau_s^{(c_i)}=q_j\},\qquad j=1,2.
\]
Conditional on $n_s$, the variables $Z_i^{(j)}$ are iid, take values in
$[-1,1]$, and satisfy
\[
    \E[Z_i^{(j)}\mid I_s]=p_s^*-r_{s,j}\geq\Delta.
\]
Hence Hoeffding's inequality implies
\[
    \P\left(\sum_{i=1}^{n_s}Z_i^{(j)}\leq0\,\middle|\,n_s\right)
    \leq \exp\left(-\frac{n_s\Delta^2}{2}\right).
\]
Let $F_s$ be the event that $\cT_s^*$ is not the unique empirical maximizer for $s$.
On $F_s$, the above event occurs for at least one $j\in\{1,2\}$. Thus, on
$n_s\geq np_{\min}/2$,
\[
    \P\left(F_s\,\middle|\,n_s\right)
    \leq 2\exp\left(-\frac{np_{\min}\Delta^2}{4}\right).
\]

If the true topology is uniquely modal for every leaf set, exact ASTRAL returns
$\cT^*$. Combining the preceding estimate with the informative-count bound
above and union bounding over the $K$ leaf sets gives
\[
    \P(\widehat{\cT}\neq\cT^*)
    \leq K\exp\left(-\frac{np_{\min}}{8}\right)
    +2K\exp\left(-\frac{np_{\min}\Delta^2}{4}\right)
    \leq 3K\exp\left(-\frac{np_{\min}\Delta^2}{8}\right),
\]
as $\Delta\leq1$. Applying the assumed lower bound on $n$
gives $\P(\widehat{\cT}=\cT^*)\geq1-\varepsilon$ as desired.
\end{proof}

\subsection{The Anomaly-Free Zone}
\label{sec:anomaly-free-zone}

Theorem~\ref{thm:correct_topo_max_expected_support} shows that $\E |M_{\cT_s^*}|>\E|M_q|$
for every $q\in A_s\setminus\{\cT_s^*\}$, but this does not imply that
%
\[\P\left(\cT_s^* = \argmax_{q' \in A_s} |M_{q'}|\right)
>\P\left(q = \argmax_{q' \in A_s} |M_{q'}|\right)\]
%
\noindent for every incorrect $q$, as it is possible $|M_{\cT_s^*}|$ is zero most of the time but takes large values
rarely. We now show that the anomaly-free property holds in the low-polymorphism regime
under the simpler $ID(\theta, \lambda, \mu; \cT^*)$ model.

Lemma~\ref{lem:quartet_supports_under_ID_process} shows in this setting that the supports take the form
%
\[|M_{uv||xy}| = N_{uv}N_{xy}, \quad |M_{ux||vy}| = N_{ux} N_{vy}, \quad |M_{uy||vx}| = N_{uy} N_{vx},\]
%
\noindent where $\{N_{uv},\ldots,N_{vx}\}$ are mutually independent Poisson random variables with means
$m_{ab}:=\E N_{ab}$. These means depend on the underlying parameters
$(\cT^*,\lambda,\mu,\theta)$ in complicated, interwoven ways. For example,
Lemma~\ref{lem:pop_sizes_in_ID_processes} gives $|I| \sim \mathrm{Poisson}(m_\theta(\ell_{ri}))$.
Hence if $\cT_s^*=uv||xy$ and it induces a balanced subtree in $\cT^*$, then
%
\[m_{uv} = \eta_{uv} \left[m_0(\ell_{ij}) + m_\theta(\ell_{ri}) \eta_{ij}(1 - \eta_{ik} (\eta_{ky} + \eta_{kx} - \eta_{kx} \eta_{ky}))\right],\]
\[m_{ux} = m_\theta(\ell_{ri})\eta_{ux} \left[(1-\eta_{jv}) (1-\eta_{ky})\right].\]
%
\noindent In the low-polymorphism regime, these means are close to zero, so the supports are heavily
concentrated on $\{0,1\}$ and the maximum support is typically the unique nonzero one.

\begin{lemma}[Low Polymorphism Implies Anomaly-Free]
\label{lem:low_polymorphism_anomaly_free}
    Suppose $\theta>0$, and consider the process
    \[Y^{(\epsilon)}\sim ID(\epsilon\theta,\epsilon\lambda,\mu;\cT^*).\]
    Fix any set of four leaves $s=\{u,v,x,y\}$.
    If
    %
    \[\omega_*^{(\epsilon)} = \E |M_{\cT_s^*}^{(\epsilon)}|, \quad \{ \omega_1^{(\epsilon)}, \omega_2^{(\epsilon)}\} = \{ \E|M_{q'}^{(\epsilon)}| \; :\; q' \in A_s \setminus \{\cT_s^*\}\},\]
    %
    \noindent are the mean supports of the correct and two incorrect topologies respectively, then for every $q\in A_s$, as $\epsilon \downarrow 0$,
    %
    \[\P\left(q = \argmax_{q' \in A_s} |M^{(\epsilon)}_{q'}| \; \Big\vert \; \max_{q' \in A_s} |M^{(\epsilon)}_{q'}| > 0\right) = \frac{\omega_q^{(1)}}{\omega_*^{(1)} + \omega_1^{(1)} + \omega_2^{(1)}} + O(\epsilon).\]
    %
    \noindent In particular, for small enough $\epsilon > 0$ the anomaly-free property~\eqref{eqn:anomaly_free_property} holds for this quartet.
\end{lemma}

\begin{proof}
    Let $N_{uv}^{(\epsilon)},\ldots,N_{vx}^{(\epsilon)}$ be the corresponding Poisson random variables from
    Lemma~\ref{lem:quartet_supports_under_ID_process}, and let
    $m_{uv}^{(\epsilon)},\ldots,m_{vx}^{(\epsilon)}$ be their means.
    By Poisson thinning, $m_{uv}^{(\epsilon)} = \epsilon m_{uv}^{(1)}$, and similarly for the other means.
    Letting $\omega_q := \omega_q^{(1)}$ and $m_{uv} := m_{uv}^{(1)}$, for $q = uv||xy$ define
    %
    \[p_{q} := \P\left(|M_{q}^{(\epsilon)}| > 0\right) = (1-e^{-\epsilon m_{uv}})(1-e^{-\epsilon m_{xy}}) = \epsilon^2 \omega_{q} + O(\epsilon^3),\]
    %
    \noindent and likewise for the other supports. More specifically, using
    $a \geq 1-e^{-a} \geq a-a^2/2$ for $a\geq0$ and
    $m_{uv}^{(1)} \leq \zeta := \theta \vee \lambda \mu^{-1}$, we obtain
    %
    \[-\epsilon^3 \zeta \omega_q \leq p_q -\epsilon^2 \omega_q \leq 0.\]
    %
    Let
    %
    \[\tau^{(\epsilon)} = \argmax_{q \in A_s} |M^{(\epsilon)}_q|, \quad I_s^{(\epsilon)} = \left\{\max_{q \in A_s} |M^{(\epsilon)}_q| > 0\right\}.\]
    %
    \noindent Since the three supports are independent, the probability that at least two
    are nonzero is at most
    %
    \[\sum_{\{q,q'\}\in\binom{A_s}{2}}p_q p_{q'}=O(\epsilon^4).\]
    %
    Consequently,
    %
    \[\P(\tau^{(\epsilon)}=q,I_s^{(\epsilon)})=p_q+O(\epsilon^4)
    =\epsilon^2\omega_q+O(\epsilon^3).\]
    %
    \noindent Likewise, with $W:=\omega_*+\omega_1+\omega_2$,
    $\P(I_s^{(\epsilon)})=\epsilon^2W+O(\epsilon^3)$. Thus
    %
    \begin{align*}
         \P\left(\tau^{(\epsilon)} = q \mid I_s^{(\epsilon)} \right)
         &=\frac{\epsilon^2\omega_q+O(\epsilon^3)}{\epsilon^2W+O(\epsilon^3)}\\[5pt]
         &=\frac{\omega_q}{\omega_*+\omega_1+\omega_2}+O(\epsilon),
    \end{align*}
    %
    \noindent as desired. Since the true topology has strictly highest expected support by
    Corollary~\ref{cor:expected_margin_under_id}, this implies the anomaly-free property for this quartet for sufficiently small $\epsilon$.
\end{proof}

Since $\cT^*$ has finitely many quartets, the anomaly-free property holds for all of them
for sufficiently small $\epsilon>0$ under the simpler
$ID(\theta, \lambda, \mu; \cT^*)$ model. This is also the regime where the immigration-death
approximation is least accurate, so it is unclear how directly the result translates to the full PCH model.
Nevertheless, it provides evidence that the anomaly-free property is reasonable, and we conjecture it
holds in the low-polymorphism regime of the full PCH model.

%%%%%%%%%%%%%%%%%%%%%%%%%%%%%%%%%%%%
% Discussion and Concluding Remarks
%%%%%%%%%%%%%%%%%%%%%%%%%%%%%%%%%%%%
\section{Discussion and Concluding Remarks}

In this chapter, we developed several potential candidates for bipartition covers which
substantially reduce the search space of the original PCH-ASTRAL inference procedure.
This modification is intended to make PCH-ASTRAL computationally feasible for much larger
collections of languages. Although restricting the search space reduces sample efficiency,
we proposed several heuristics designed to mitigate this loss in performance.

Beyond this, we derived a sample-complexity bound for PCH-ASTRAL under both the PCH model
and its immigration-death approximation. We also obtained a bound for the character-voting
procedure PCH-ASTRAL(sc) throughout the anomaly-free zone and verified this condition under
the ID model in the low-polymorphism regime. In both cases, the rates depend explicitly on a
quartet-level margin measuring the advantage of the correct quartet topology over the
incorrect ones. We further gave partial bounds relating these margins to the underlying
model parameters.

There are several natural directions for future work. First, the PCH tail parameter
$(1+\||c(r)|\|_{\Psi_1}+\lambda/\mu)^2$ from
Lemma~\ref{lem:pch_subweibull_quartet_supports} is coarse. It would be useful to sharpen it to retain
pair-specific survival factors, as the exact Poisson representation does under the ID approximation.

Second, since the true tree $\cT^*$ is unknown at inference time, it would be useful
to develop topology-free sample-complexity bounds depending only on simple geometric
 parameters of $\cT^*$, such as the minimum branch length, analogous to the results of
 \cite{shekhar2018_how_many_genes_are_enough}. Under the ID approximation, the Poisson
 means $m_{uv}=\E N_{uv}$ are explicit functions of the model parameters, so more careful
bookkeeping in Lemma~\ref{lem:sample_complexity_pch_astral} may lead to relatively explicit
bounds.

Third, the anomaly-free condition for PCH-ASTRAL(sc) remains somewhat opaque. Under the ID
approximation, the condition reduces to a comparison of independent product-Poisson random
variables, and is therefore numerically checkable from the quartet-support means. However,
we do not currently have simple sufficient conditions ensuring that the anomaly-free
 property holds.

Extensive empirical simulations would also be valuable, both to compare PCH-ASTRAL and PCH-ASTRAL(sc) under various
model regimes and to assess how well the theoretical margins predict finite-sample behavior.
It would also be interesting to study how well our proposed bipartition heuristics improve
coverage properties.
